\documentclass[12pt]{unisenaiclass}
\usepackage{unisenaistyle}
\usepackage[utf8]{inputenc}
\usepackage[brazil]{babel}
\usepackage{graphicx}
\usepackage{float}
\usepackage{amsmath}
\usepackage{hyperref}
\usepackage{caption}
\usepackage{csquotes}
\usepackage[backend=biber, style=abnt]{biblatex}
\addbibresource{referencias.bib}

\title{DESENVOLVIMENTO DE UM HUB IoT COM BACKEND NODE.JS, PROTOCOLO OPC-UA E INTERFACE WEB NEXT.JS}
\subtitle{DEVELOPMENT OF AN IoT HUB WITH NODE.JS BACKEND, OPC-UA PROTOCOL AND NEXT.JS WEB INTERFACE}

\author{
  \textbf{Lauren Julia Reis Silva} \thanks{Estudante do curso de Tecnologia, SENAI.}
}

\begin{document}
\maketitle

% ─────────────────────────────────────────────────────────────
%  RESUMO
% ─────────────────────────────────────────────────────────────
\begin{resumo}
Este artigo apresenta o desenvolvimento de um hub IoT integrado a um backend Node.js e
uma interface web construída com Next.js. O sistema permite o monitoramento em tempo real
de sensores de temperatura, pressão, umidade e presença, além do controle remoto de um
atuador (trava de segurança). A comunicação entre os dispositivos físicos e a plataforma
digital é intermediada pelo protocolo OPC-UA, com um servidor Python expondo variáveis
ao Node-RED, que encaminha os dados à API RESTful via requisições HTTP. A metodologia
envolveu a criação de uma API RESTful com Express.js, a conteinerização completa dos
serviços com Docker Compose e a construção de uma interface responsiva com React e Tailwind
CSS. Os resultados demonstram a viabilidade técnica da solução para aplicações de automação
residencial, industrial e agrícola, evidenciando escalabilidade e interoperabilidade entre
tecnologias heterogêneas.
\end{resumo}

\begin{justify}
\textbf{Palavras-chave:} IoT; Backend; API; OPC-UA; Node-RED; Docker; Sensores; Automação.
\end{justify}

% ─────────────────────────────────────────────────────────────
%  ABSTRACT
% ─────────────────────────────────────────────────────────────
\begin{abstract1}
This paper presents the development of an IoT hub integrated with a Node.js backend and a
Next.js web interface. The system enables real-time monitoring of temperature, humidity,
pressure, and presence sensors, as well as remote control of a security lock actuator.
Communication between physical devices and the digital platform is mediated by the OPC-UA
protocol, with a Python server exposing variables to Node-RED, which forwards data to a
RESTful API via HTTP requests. The methodology involved the creation of a RESTful API using
Express.js, full containerization of services with Docker Compose, and the construction of
a responsive interface with React and Tailwind CSS. Results demonstrate the technical
feasibility of the solution for residential, industrial, and agricultural automation
applications, highlighting scalability and interoperability across heterogeneous technologies.
\end{abstract1}

\begin{justify}
\textbf{Keywords:} IoT; Backend; API; OPC-UA; Node-RED; Docker; Sensors; Automation.
\end{justify}

% ─────────────────────────────────────────────────────────────
%  1. INTRODUÇÃO
% ─────────────────────────────────────────────────────────────
\section{Introdução}

A Internet das Coisas (IoT) tem se consolidado como uma tecnologia essencial para a
automação e o monitoramento inteligente de ambientes industriais, agrícolas e residenciais.
A capacidade de conectar dispositivos físicos à internet e processar seus dados em tempo
real representa um avanço significativo na forma como sistemas de controle são projetados
e operados.

Este trabalho apresenta o desenvolvimento de um sistema IoT completo, composto por quatro
camadas integradas: um servidor OPC-UA em Python responsável pela simulação e exposição
de dados de sensores; o Node-RED como middleware de integração; uma API RESTful desenvolvida
em Node.js com Express; e uma interface web construída com Next.js e React. O objetivo
central é demonstrar a viabilidade técnica de integrar tecnologias abertas e amplamente
adotadas para construir uma plataforma de monitoramento e controle escalável.

O protocolo OPC-UA (\textit{Open Platform Communications Unified Architecture}) foi
escolhido por ser o padrão industrial para comunicação segura e interoperável entre
dispositivos de campo e sistemas de supervisão \cite{aguiar2023}. A arquitetura proposta
permite que sensores de temperatura, pressão, umidade e presença tenham seus dados
coletados, armazenados e visualizados em um painel \textit{web} de forma contínua.

% ─────────────────────────────────────────────────────────────
%  2. REVISÃO DE LITERATURA
% ─────────────────────────────────────────────────────────────
\section{Revisão de Literatura}

Segundo \textcite{aguiar2023}, a Internet das Coisas pode ser descrita como uma rede de
objetos físicos conectados à internet, capazes de interagir por meio de sensores e
softwares. O autor destaca que a utilização de APIs RESTful é fundamental para conectar o
mundo físico ao digital, permitindo que dispositivos troquem dados de forma eficiente e
escalável.

De acordo com \textcite{rodrigues2024}, a expansão de funcionalidades em sistemas IoT
depende diretamente da integração com APIs HTTP, que possibilitam maior flexibilidade e
controle remoto dos dispositivos. A autora enfatiza que o uso de arquiteturas REST garante
interoperabilidade entre diferentes plataformas e tecnologias, tornando os sistemas mais
adaptáveis às necessidades dos usuários.

A literatura recente sobre IoT aponta que a integração com APIs RESTful é essencial para
garantir interoperabilidade e escalabilidade em ambientes complexos, permitindo que sensores
e atuadores sejam facilmente conectados a aplicações web e móveis \cite{aguiar2023,
rodrigues2024}.

No que tange ao protocolo OPC-UA, trata-se de um padrão amplamente adotado na indústria
4.0 por oferecer comunicação orientada a serviços com suporte a segurança, semântica de
dados e independência de plataforma. A combinação desse protocolo com ferramentas de
orquestração como o Node-RED e a conteinerização via Docker representa uma abordagem
moderna para sistemas de supervisão e controle.

\begin{itemize}
    \item \textbf{Exemplo de citação direta curta:}
\end{itemize}

Segundo \textcite{rodrigues2024}, \enquote{a interoperabilidade em IoT só é possível por
meio de APIs RESTful bem definidas}.

\begin{quote}
\textbf{Citação direta longa:} \textcite{rodrigues2024} afirma:

\enquote{A utilização de APIs RESTful em sistemas IoT não apenas garante a
interoperabilidade entre dispositivos, mas também possibilita a escalabilidade e a
segurança necessárias para aplicações críticas. Sem esse padrão de comunicação, os sistemas
tornam-se fragmentados e pouco confiáveis.}

\textbf{Citação indireta:} De acordo com \textcite{rodrigues2024}, a integração de APIs
RESTful em sistemas IoT é fundamental para assegurar interoperabilidade e escalabilidade,
permitindo que diferentes dispositivos troquem informações de maneira segura e eficiente.
\end{quote}

O sistema TeX foi desenvolvido por Knuth com o objetivo de produzir livros com alta
qualidade tipográfica \cite{knuth1984tex}. Ao final do documento, as referências serão
listadas automaticamente.

% ─────────────────────────────────────────────────────────────
%  3. METODOLOGIA
% ─────────────────────────────────────────────────────────────
\section{Metodologia}

O sistema foi desenvolvido em quatro camadas integradas, conforme descrito a seguir.

\subsection{Servidor OPC-UA (Python)}

O servidor de dados foi implementado em Python utilizando a biblioteca \texttt{opcua}
(versão 0.98.13). O servidor expõe cinco variáveis no namespace
\texttt{http://ads.freeopcua.server}, mapeadas ao objeto \texttt{Sensor}:
\texttt{temperatura} (Float), \texttt{pressao} (Float), \texttt{umidade} (Float),
\texttt{sensor\_presenca} (Boolean) e \texttt{trava\_seguranca} (Boolean). Os valores são
gerados de forma aleatória a cada segundo — temperatura entre 18\,°C e 30\,°C, pressão
entre 0,8 e 1,5 bar, umidade entre 35\% e 60\% — simulando um ambiente industrial real.
O endpoint do servidor é \texttt{opc.tcp://0.0.0.0:4840}, sem camada de segurança, para
facilitar a integração em ambiente de desenvolvimento.

\subsection{Node-RED (Middleware de Integração)}

O Node-RED atua como intermediário entre o servidor OPC-UA e a API REST. Fluxos
configurados leem os valores de temperatura e pressão do servidor OPC-UA a cada 2 segundos,
formatam os dados em JSON e os encaminham ao endpoint \texttt{POST /newData} da API via
nó \textit{HTTP Request}. Isso desacopla o protocolo industrial OPC-UA do protocolo web
HTTP, permitindo que ambas as camadas evoluam de forma independente.

\subsection{Backend — API RESTful (Node.js / Express)}

A API foi desenvolvida com Node.js e Express.js (v5.2.1), exposta na porta 8080. Ela
recebe, armazena em memória e fornece os dados dos sensores, além de gerenciar um cadastro
de dispositivos IoT. As rotas implementadas estão descritas na Tabela \ref{tab:rotas}.

\begin{table}[H]
\centering
\caption{Rotas da API RESTful}
\label{tab:rotas}
\begin{tabular}{|c|c|l|}
\hline
\textbf{Método} & \textbf{Rota}              & \textbf{Descrição} \\ \hline
GET    & \texttt{/iot}               & Lista todos os dados de sensores armazenados \\ \hline
GET    & \texttt{/sensor/:id}        & Retorna dado de sensor por índice \\ \hline
POST   & \texttt{/newData}           & Recebe dados do Node-RED e armazena \\ \hline
PUT    & \texttt{/sensor/:id}        & Atualiza dado de sensor por ID \\ \hline
GET    & \texttt{/devices}           & Lista dispositivos cadastrados \\ \hline
POST   & \texttt{/devices}           & Cadastra novo dispositivo \\ \hline
PUT    & \texttt{/devices/:id}       & Atualiza dados de um dispositivo \\ \hline
PATCH  & \texttt{/devices/:id/status}& Altera o status do dispositivo \\ \hline
DELETE & \texttt{/devices/:id}       & Remove um dispositivo \\ \hline
\end{tabular}
\end{table}

O body esperado no endpoint \texttt{POST /newData} segue o seguinte formato JSON:

\begin{verbatim}
{
  "temperatura": 25.4,
  "pressao": 1.2,
  "umidade": 55.0,
  "sensor_presenca": true,
  "trava_seguranca": false
}
\end{verbatim}

\subsection{Frontend — Interface Web (Next.js / React)}

A interface foi desenvolvida com Next.js 16 e React 19, utilizando Tailwind CSS para
estilização. A aplicação é composta por três módulos principais:

\begin{itemize}
    \item \textbf{Módulo de Usuários (\texttt{/})}: formulário de cadastro e listagem de
    usuários com modal de edição.

    \item \textbf{Módulo de Dispositivos (\texttt{/devices})}: CRUD completo de
    dispositivos IoT. Permite cadastrar sensores e atuadores com nome, tipo, endereço IP
    ou NodeID OPC-UA, e descrição. O status de cada dispositivo pode ser alternado entre
    \textit{online}, \textit{offline} e \textit{alerta}.

    \item \textbf{Dashboard de Sensores (\texttt{/dashboard})}: exibe em tempo real as
    últimas leituras de temperatura, pressão, umidade, sensor de presença e trava de
    segurança. Os dados são atualizados via \textit{polling} a cada 2 segundos utilizando
    \texttt{useEffect} e \texttt{setInterval}. Também apresenta um histórico completo de
    leituras em tabela.
\end{itemize}

\subsection{Conteinerização com Docker Compose}

Todos os serviços são orquestrados por um arquivo \texttt{docker-compose.yaml}, que define
quatro contêineres em uma rede interna \texttt{ifaci-net}: \texttt{opcua-server} (porta
4840), \texttt{api-ifaci} (porta 8080), \texttt{nodered-ifaci} (porta 1880) e
\texttt{frontend-ifaci} (porta 3000). O uso de Docker garante portabilidade e reprodutibilidade
do ambiente, eliminando dependências de configuração manual.

% ─────────────────────────────────────────────────────────────
%  4. RESULTADOS E DISCUSSÕES
% ─────────────────────────────────────────────────────────────
\section{Resultados e Discussões}

O sistema integrado apresentou os seguintes resultados após testes de funcionamento:

\begin{itemize}
    \item \textbf{Monitoramento em tempo real}: o painel \textit{dashboard} exibiu os dados
    dos sensores com latência inferior a 2 segundos, demonstrando a eficácia do
    \textit{polling} periódico via \texttt{fetch}.

    \item \textbf{Gerenciamento de dispositivos}: o módulo CRUD de dispositivos permitiu
    cadastro, edição, remoção e alternância de status (online/offline/alerta) com
    confirmação visual imediata na interface.

    \item \textbf{Integração OPC-UA $\rightarrow$ Node-RED $\rightarrow$ API}: os dados
    simulados pelo servidor Python foram corretamente lidos pelo Node-RED e encaminhados
    à API, comprovando a viabilidade do pipeline de dados proposto.

    \item \textbf{Conteinerização}: todos os serviços subiram sem conflitos a partir do
    comando \texttt{docker compose up}, validando a portabilidade da solução.
\end{itemize}

\subsection{Exemplo de leitura dos sensores}

A Tabela \ref{tab:sensores} apresenta valores representativos coletados durante os testes
do sistema, com os cinco sensores monitorados simultaneamente.

\begin{table}[H]
\centering
\caption{Monitoramento de temperatura, umidade, pressão, presença e trava de segurança.}
\label{tab:sensores}
\begin{tabular}{|c|c|c|}
\hline
\textbf{Sensor}     & \textbf{Valor Atual} & \textbf{Unidade} \\ \hline
Temperatura         & 25,40                & °C               \\ \hline
Pressão             & 1,20                 & bar              \\ \hline
Umidade             & 55,00                & \%               \\ \hline
Sensor de Presença  & Detectada            & —                \\ \hline
Trava de Segurança  & Livre                & —                \\ \hline
\end{tabular}
\end{table}

Os resultados confirmam que a arquitetura em camadas adotada — com OPC-UA para comunicação
industrial, Node-RED como middleware, Express como backend e Next.js como frontend —
oferece boa separação de responsabilidades e facilita a manutenção e expansão do sistema.
A utilização de Docker Compose simplificou significativamente o processo de implantação,
alinhando-se às práticas modernas de desenvolvimento orientado a microsserviços
\cite{rodrigues2024}.

% ─────────────────────────────────────────────────────────────
%  5. CONCLUSÃO
% ─────────────────────────────────────────────────────────────
\section{Conclusão}

O projeto demonstrou a viabilidade técnica de integrar sensores e atuadores industriais em
uma plataforma IoT completa, combinando o protocolo OPC-UA, o Node-RED, uma API RESTful em
Node.js e uma interface web moderna em Next.js. A solução atende aos requisitos de
monitoramento em tempo real, gerenciamento de dispositivos e controle remoto de atuadores,
sendo conteinerizada integralmente com Docker.

Como trabalhos futuros, pretende-se expandir a solução para incluir: persistência de dados
em banco de dados relacional (PostgreSQL) ou de séries temporais (InfluxDB); análise
preditiva com modelos de \textit{machine learning} para detecção de anomalias nos sensores;
integração com serviços em nuvem como AWS IoT Core ou Azure IoT Hub; e autenticação de
usuários com JWT para controle de acesso por níveis de permissão.

% ─────────────────────────────────────────────────────────────
%  AGRADECIMENTOS
% ─────────────────────────────────────────────────────────────
\section*{Agradecimentos}

Agradeço à instituição de ensino SENAI e ao orientador Gabriel, que contribuiu de forma
fundamental para o planejamento e desenvolvimento deste trabalho.

% ─────────────────────────────────────────────────────────────
%  REFERÊNCIAS
% ─────────────────────────────────────────────────────────────
\printbibliography

\end{document}
